% part: intuitionistic-logic
% chapter: propositions-as-types
% section: sequent-natural-deduction

\documentclass[../../../include/open-logic-section]{subfiles}

\begin{document}

\olfileid{int}{pty}{snd}

\olsection{সিকোয়েন্টে স্বাভাবিক নিষ্পাদন}

$\Gamma \Sequent !A$ লিখি যদি এমন একটি স্বাভাবিক
নিষ্পাদনের !!{derivation} থাকে যার !!{undischarged}
অনুমানগুলির সমষ্টি $\Gamma$ এবং সিদ্ধান্ত $!A$;
$\Gamma$ ফাঁকা হলে লিখি $\Sequent !A$।

$\Gamma \cup \{!A_1, \dots, !A_n\}$-কে
$\Gamma, !A_1, \dots, !A_n$ লিখি; আর
$\Gamma \cup \Delta$-কে লিখি $\Gamma, \Delta$।

ধরি, $\Gamma \Sequent !A \land !A$ আছে, অর্থাৎ
$\Gamma$-কে !!{undischarged} অনুমানের সমষ্টি এবং
$!A \land !A$-কে অন্তিম-!!{formula} রেখে
!!a{derivation} আছে। তার শেষে $\Elim{\land}$
প্রয়োগ করলে একই !!{undischarged} অনুমান থেকে
$!A$ সিদ্ধান্তযুক্ত !!a{derivation} পাই। অন্য
কথায়, $\Gamma \Sequent !A \land !B$ থেকে
$\Gamma \Sequent !A$ পাওয়া যায়।
\begin{prooftree}
  \Axiom$\Gamma \fCenter !A \land !B$
  \RightLabel{$\Elim{\land}$}
  \UnaryInf$\Gamma \fCenter !A$
  \DisplayProof\qquad\bottomAlignProof
  \Axiom$\Gamma \fCenter !A \land !B$
  \RightLabel{$\Elim{\land}$}
  \UnaryInf$\Gamma \fCenter !B$
\end{prooftree}
% BN-SRC-619: optional earlier alteration withdrawn; source restored.
$\Elim{\land}$ লেবেলটি স্বাভাবিক নিষ্পাদনের
একই নামের বিধির সঙ্গে এই বিধির সম্পর্ক দেখায়।

একইভাবে, ধরি $\Gamma, !A \Sequent !B$ আছে;
অর্থাৎ $\Gamma, !A$ হলো !!{undischarged}
অনুমানের সমষ্টি এবং $!B$ হলো অন্তিম-!!{formula}
এমন !!a{derivation} আছে। $\Intro{\lif}$ বিধি
প্রয়োগ করলে $\Gamma$-কে !!{undischarged} অনুমান
এবং $!A \lif !B$-কে অন্তিম-!!{formula} রেখে
!!a{derivation} পাই, অর্থাৎ $\Gamma \Sequent
!A \lif !B$। এই লিখনরীতিতে অনুমানের
!!{discharge} আরও স্পষ্ট হয়েছে।
\begin{prooftree}
  \Axiom $\Gamma, !A \fCenter !B$
  \RightLabel{$\Intro{\lif}$}
  \UnaryInf $\Gamma \fCenter !A \lif !B$
\end{prooftree}

অন্য বিধি থেকেও একইভাবে সিদ্ধান্ত পাওয়া যায়;
সবগুলি নিচে দেখানো হলো:
\begin{gather*}
  \Axiom $\Gamma \fCenter !A$
  \Axiom $\Delta \fCenter !B$
  \RightLabel{$\Intro{\land}$}
  \BinaryInf $\Gamma,\Delta \fCenter !A \land !B$
  \DisplayProof\\
  \Axiom $\Gamma \fCenter !A \land !B$
  \RightLabel{$\Elim{\land}_1$}
  \UnaryInf $\Gamma \fCenter !A$
  \DisplayProof
  \qquad
  \Axiom $\Gamma \fCenter !A \land !B$
  \RightLabel{$\Elim{\land}_2$}
  \UnaryInf $\Gamma \fCenter !B$
  \DisplayProof
  \\
  \Axiom $\Gamma \fCenter !A$
  \RightLabel{$\Intro{\lor}_1$}
  \UnaryInf $\Gamma \fCenter !A \lor !B$
  \DisplayProof\qquad
  \Axiom $\Gamma \fCenter !B$
  \RightLabel{$\Intro{\lor}_2$}
  \UnaryInf $\Gamma \fCenter !A \lor !B$
  \DisplayProof\\
  \Axiom $\Gamma \fCenter !A \lor !B$
  \Axiom $\Delta, !A \fCenter !C$
  \Axiom $\Delta', !B \fCenter !C$
  \RightLabel{$\Elim{\lor}$}
  \TrinaryInf $\Gamma, \Delta, \Delta' \fCenter !C$
  \DisplayProof
  \\
  \Axiom $\Gamma, !A \fCenter !B$
  \RightLabel{$\Intro{\lif}$}
  \UnaryInf $\Gamma \fCenter !A \lif !B$
  \DisplayProof
  \qquad
  \Axiom $\Delta \fCenter !A \lif !B$
  \Axiom $\Gamma \fCenter !A$
  \RightLabel{$\Elim{\lif}$}
  \BinaryInf $\Gamma, \Delta \fCenter !B$
  \DisplayProof
  \\
  \Axiom $\Gamma \fCenter \lfalse$
  \RightLabel{$\FalseInt$}
  \UnaryInf $\Gamma \fCenter !A$
  \DisplayProof
\end{gather*}

যেকোনো অনুমান নিজেই $!A$ থেকে $!A$-এর
!!a{derivation}; তাই সব সময় $!A \Sequent !A$।

\begin{prooftree}
  \AxiomC{$ $}
  \UnaryInfC{$!A \fCenter !A$}
\end{prooftree}

বিধিগুলিকে একত্রে কোন কোন স্বাভাবিক নিষ্পাদনের
!!{derivation}s আছে তার একটি কলন হিসেবে ধরা
যায়। এগুলি স্বাভাবিক নিষ্পাদনেরই বিকল্প
লিখনরীতি: প্রতিটি ধাপে কেবল যে
!!{formula} !!{derive}d হয়েছে সেটিই নয়, যে !!{undischarged} অনুমানগুলি
থেকে সেটি !!{derive} করা হয়েছে, সেগুলিও লেখা
থাকে।

\begin{prooftree}
  \Axiom$!A \fCenter !A$
  \UnaryInf $!A \fCenter !A \lor (!A \lif \lfalse)$
  \Axiom $!B \fCenter !B$
  \BinaryInf$!A, !B \fCenter \lfalse$
  \UnaryInf$!B \fCenter !A \lif \lfalse$
  \UnaryInf$!B \fCenter !A \lor (!A \lif \lfalse)$
  \Axiom$!B \fCenter !B$
  \BinaryInf$!B \fCenter \lfalse$
  \UnaryInf$\fCenter !B \lif \lfalse$
\end{prooftree}
% BN-SRC-620 TARGET-CORRECTION-BEGIN
\par\smallskip\noindent\textnormal{\small\textbf{উৎসসংশোধন (BN-SRC-620):} শেষ প্রমাণবৃক্ষে উৎসের অসম্পূর্ণ বন্ধনী, ভুল প্রসঙ্গ \( !B\lif \), চিহ্নগুলি সংশোধন করা হয়েছে। \( !B=( !A\lor(!A\lif\lfalse))\lif\lfalse \) ধরে প্রথমে \( !A,!B\Sequent\lfalse \), তারপর \( !B\Sequent !A\lif\lfalse \), পরে \( !B\Sequent !A\lor(!A\lif\lfalse) \), শেষে \( \Sequent !B\lif\lfalse \) পাওয়া যায়।}
% BN-SRC-620 TARGET-CORRECTION-END
এখানে $!B$ হলো $(!A \lor (!A \lif \lfalse))\lif \lfalse$-এর সংক্ষিপ্ত রূপ।

\end{document}
